%%%%%%%%%%%%%%%%%%%%%%%%%%%%%%%%%%%%%%%%%
% Author: bmarron
% Created: 02 Nov 2026
% Modified:
%
% 12-point; 1in margins
% One inch vertical space is 72.27pt
%
%%%%%%%%%%%%%%%%%%%%%%%%%%%%%%%%%%%%%%%%%



%----------------------------------------------------------------------------------------
%	PACKAGES AND OTHER DOCUMENT CONFIGURATIONS
%----------------------------------------------------------------------------------------

\documentclass[12pt, letterpaper]{article}

\usepackage[left=1in,right=1in,top=1in,bottom=1in]{geometry}       % Document margins

\usepackage{afterpage}
%\usepackage{datetime2}               % specialty date formats
\usepackage{lipsum}

%BIBLIOGRAPHY (new way; biber+biblatex)
\usepackage[utf8]{inputenc}
\usepackage[T1]{fontenc}              % Use 8-bit encoding that has 256 glyphs
\usepackage[english]{babel}
\usepackage[style=apa,			          % <== APA style for documents
            sorting=nyt,
            backref=true,           
            doi=false,
            isbn=false,
            url=false,
            backend=biber]{biblatex}
\addbibresource{~/Desktop/GenRepo_01/BibLatex/MyCurrent_Library/My_Library_20201129.bib} 
\DeclareLanguageMapping{english}{american-apa}


\usepackage{lineno}
\usepackage{csquotes}

\usepackage{amsmath}
\usepackage{amsthm}
\usepackage{amssymb}
\usepackage{pdfpages}
\usepackage{verbatim}

\usepackage{outlines}


\usepackage{threeparttable}
\usepackage{pgfgantt}                 % simple gantt charts
\usepackage{bigfoot}
\usepackage{multirow}
\usepackage{microtype}                % Slightly tweak font spacing for aesthetics
\usepackage{booktabs}                 % Horizontal rules in tables
\usepackage{float}                    % specific locations with the [H] (e.g. \begin{table}[H])
\usepackage{paralist}                 % bullet points with less space between them
\usepackage[hang, 
            small,
            labelfont=bf,
            up,
            textfont=it,up]{caption}    % Custom captions under/above floats in tables or figures

\usepackage[sc]{mathpazo}             % Use the Palatino font and the Pazo fonts for math
\linespread{1.05}                     % global setting - Palatino needs more space
\usepackage{setspace}
%\doublespacing                       % global setting

\usepackage{abstract}                                               % Allows abstract customization
\renewcommand{\abstractnamefont}{\normalfont\bfseries}              % Set the "Abstract" text to bold
\renewcommand{\abstracttextfont}{\normalfont\small\itshape}         % Set the abstract itself to small italic text

\usepackage{titlesec}                                               % Allows customization of titles
\renewcommand\thesection{\Roman{section}}                           % Roman numerals for the sections
\renewcommand\thesubsection{\arabic{subsection}}                    % Arabic numerals for subsections
\titleformat{\section}[block]{\large\scshape\centering}{\thesection.}{1em}{}        % Change the look of the section titles
\titleformat{\subsection}[block]{\large\underline}{\thesubsection}{1em}{}                    % Change the look of the subsection titles
\titlespacing*{\subsubsection}{0pt}{0.75\baselineskip}{0.25\baselineskip}            % Change spacing after section


\usepackage{fancyhdr}                                              % Headers and footers
\pagestyle{fancy}                                                 % All pages have headers and footers
%\pagestyle{plain}                                                  % globally set to "plain"
%\fancyhead{} % Blank out the default header
%\fancyfoot{} % Blank out the default footer
\fancyhead[L]{Applied Remote Sensing for Agroecology}                                                                                    
\fancyhead[R]{\includegraphics[scale= 0.2]{graphics/UMA_Logo.png}}
%\fancyhead[C]{}                                                   % Custom header text
\fancyfoot[C]{\thepage}    
%\fancyfoot[RO,LE]{\thepage}                                        % Custom footer text



\usepackage{everypage}                             % Required for watermarks
\usepackage{draftwatermark}                        % Watermarks or not
\SetWatermarkLightness{0.95}
\SetWatermarkScale{0.75}                            % Set to 1.0 for default
%\SetWatermarkText{FINAL DRAFT}                      % comment out for default (DRAFT)   

%\usepackage{multicol}                                      % Used for the two-column layout of the document
%\usepackage{hyperref}                                      % Interferes with cite.sty 




%-------------------------------------------------------------
% NEW COMMANDS
%---------------------------------------------------------------------
%This command creates a box marked ``To Do'' around text.
%To use: 
% type \todo{  insert text here  }.
\newcommand{\todo}[1]{\vspace{5 mm}\par \noindent
\marginpar{\textsc{To Do}}
\framebox{\begin{minipage}[c]{0.95 \textwidth}
\tt\begin{center} #1 \end{center}\end{minipage}}\vspace{5 mm}\par}

% This command makes a sun symbol
% To use:
% Type \sun
\newcommand{\sun}{\ensuremath{\odot}}                                   
\newcommand\ddfrac[2]{\frac{\displaystyle #1}{\displaystyle #2}}

% This command creates a custom indent
% To use:
%\begin{myindentpar}{2em}
% blah blah blah
%\end {myindentpar}
\newenvironment{myindentpar}[1]%
   {\begin{list}{}%
       {\setlength{\leftmargin}{#1}}%
           \item[]%
   }
     {\end{list}}

%----------------------------------------------------------------------------------------
%	TITLE SECTION
%----------------------------------------------------------------------------------------

\title{\vspace{-5mm}\fontsize{14pt}{12pt}
                    \selectfont \textbf{Introduction to Remote Sensing in Agriculture}\\  % Article title
                   \vspace{2mm}
                   Bruce D. Marron \\
                   \vspace{2mm}
                   \textit{(Course Outline P01)}
      } 

\author{
%    \textit{(Course Outline P01)} \\                     % Name in “small caps”
%     \normalsize FOR 803-001 \\
%     \vspace{-9mm}
%      \normalsize North Carolina State University      % Institution or Class ID
     \vspace{-20mm}
}
\date{}



%---------------------------------------------------------------------------------------
% HELPFUL
%----------------------------------------------------------------------------------------

%\enquote*{text}                                 % single quotes
%\shortcites{citation}
%\cite{citation}                                % gives Authors et al.

%\includegraphics[width=6.0in, height=5in, page=2]{graphics/Problems1.pdf}              % pdf
%\includegraphics[width=0.35\textwidth, height=0.6\textheight]{graphics/fig1}           % pdf
%\thispagestyle{plain}                          % headers/footers included page or not (default =TRUE except for 1st)



%-------- TEMPLATE -----------------------------
%%%%%%%%%%%%%%%%%%%%%%%%%
%\section{}
%%%%%%%%%%%%%%%%%%%%%%%%%
%-------------------------------------------
%\large
%\begin{flushleft}
%\textit{}
%\end{flushleft}
%\normalsize
%--------------------------------------------

%\begin{outline}[enumerate]
%\1 \textbf{}

%\end{outline}




%========================
% DOCUMENT
%=========================
\begin{document}


\begin{singlespacing}
\maketitle                                      % insert title or not
\end{singlespacing}
\thispagestyle{fancy}                            % headers/footers included or not (default=FALSE for first pg)
%\linenumbers                                    % line numbers for document or not

%---------- DOC TEXT -----------------------

%####################################
\section{Theme: Agriculture in the 21st C.}
%#####################################

%---------------------------------------
\large
\begin{flushleft}
\textbf{\textit{Section: Ag's Challenges and Solutions}}
\end{flushleft}
\normalsize
%---------------------------------------

\begin{outline}[enumerate]
\1 Challenges
\2
\1 Solutions
\2 


\end{outline}


%---------------------------------------
\large
\begin{flushleft}
\textbf{\textit{Section: Introduction to Remote Sensing in Ag}}
\end{flushleft}
\normalsize
%---------------------------------------

\begin{outline}[enumerate]

\1 Satellites, Satellite Data, and Satellite Maps
\2 Satellites
\3 Sentinel-2
\3 Landsat 8
\2 Satellite Data
\3 Applied Geospatial Analysis from Satellite Data
\2 Satellite Maps
\3 

\1 Drones, Drone Data, and Drone Maps
\2 Drones
\3
\2 Drone Data
\3 
\2 Drone Maps
\3 Applied Geospatial Analysis from Drone Data
\2 
\end{outline}

%####################################
\section{Theme: AI in the 21st C.}
%####################################

%-------------------------------------
\large
\begin{flushleft}
\textbf{\textit{Section: Introduction to Professional AI Use}}
\end{flushleft}
\normalsize
%---------------------------------------

\begin{outline}[enumerate]
\1 AI Providers: Models and APIs
\2 

\end{outline}


%####################################
\section{Theme: Records Management in the 21st C.}
%####################################
%-------------------------------------------
\large
\begin{flushleft}
\textbf{\textit{Introduction to Records and Document Management}}
\end{flushleft}
\normalsize
%--------------------------------------------

\begin{outline}[enumerate]
\1 QAQC for Docs
\1 GitHub

\end{outline}

\pagebreak

%----------------------------------------------------------------------------------------
%	NOTES
%----------------------------------------------------------------------------------------
\thispagestyle{plain} 

\large
\begin{flushleft}
\textit{NOTES}
\end{flushleft}
\normalsize

\begin{itemize}
  \item remote sensing using multispectral images is a proxy for extensive manual crop monitoring operations in the field and provides potential savings in precious time and resources. Precision agriculture has grown and established itself as a central tool in agricultural management to forecast yield, monitor plant stress, and optimize fertilization, irrigation, and soil tilling activities for a more sustainable management of farming practices and improvements in yield for farmers.(\textbf{\cite{segarra_remote_2020}})
  \item Sentinel-2 platform is the pivotal sensor: by providing freely accessible multispectral imagery with 10–20 m spatial resolution, frequent revisit times, and red-edge bands, Sentinel-2 has enabled a shift from regional assessments toward field-level and within-field yield monitoring, and its red-edge bands are also effective for monitoring nitrogen status. The Sentinel-2 constellation, with improved spatial, spectral, and temporal resolution was designed specifically to meet the needs of the agricultural community, both farmers and academic researchers with a focus on international agricultural development. (\textbf{\cite{segarra_remote_2020}})
  \item One very notable use of Sentinel-2 features at field monitoring is the ESA Sen2-Agri system which generates modules of time series along the growing season. These different products consist of a monthly cloud-free composite of surface reflectance at 10–20 m resolution; monthly dynamic cropland masks, delivered from the agricultural mid-season onwards; a cultivated crop type map at 10 m resolution for main crop groups; and time series of vegetation status indicators, NDVI and LAI (\textbf{\cite{segarra_remote_2020}})
  \item Concerning agriculture, in situations where crops interact with any given aspect of their environment (seasonal climatic variations, meteorological extreme events, pests, soil properties, etc.) or as crops grow and pass through different phenological stages, the interactions between plants and light reflectance translate into changes in plant signal patterns that can be interpreted using satellite data.(\textbf{\cite{segarra_remote_2020}})
\end{itemize}




%----------------------------------------------------------------------------------------
%	REFERENCE LIST
%----------------------------------------------------------------------------------------
\clearpage
\thispagestyle{plain} 
%\renewcommand{\thepage}{}
\printbibliography


\end{document}
